% 编译器：XeLaTeX。以下字号、页边距、章节为可改默认值。
\documentclass[UTF8,a4paper,zihao=-4]{ctexart}
\usepackage[margin=2.5cm]{geometry}
\usepackage{amsmath,amssymb,booktabs,graphicx,caption,float,setspace,ragged2e}
\setstretch{1.3}
\captionsetup[figure]{position=bottom}
\captionsetup[table]{position=top}
\newcommand{\pending}[1]{\textbf{【待补：#1】}}
\title{\pending{实验名称}实验报告}
\author{\pending{姓名、学号}}
\date{\pending{实验日期}}
\begin{document}
\maketitle
\justifying
\begin{abstract}
\justifying
\pending{目的、方法、真实核心结果和结论。无数据时不填写测量结果。}
\end{abstract}
\section{实验目的与原理}
\pending{物理关系、适用条件、必要公式及符号定义。用 equation 环境居中并自动编号，用 label 和 eqref 引用。}
\section{仪器与方法}
\pending{实际仪器、精度、实验条件和操作。}
\section{数据与结果}
\begin{table}[htbp]
\centering
\caption{原始测量记录（待补真实数据）}
\label{tab:raw}
\begin{tabular}{ccc}
\toprule
序号 & \pending{物理量/单位} & \pending{物理量/单位} \\
\midrule
\pending{记录} & \pending{数据} & \pending{数据} \\
\bottomrule
\end{tabular}
\end{table}
\pending{说明数据处理方法、结果及有依据的不确定度。}
% 使用 figure 环境时先 includegraphics 后 caption，配合 [H] 等控制图注同页。
% 坐标轴必须含量和单位。没有实际数据图时不要插入占位图片。
\section{讨论与思考题}
\pending{根据实际数据讨论误差和改进，并回答讲义问题。}
\section{结论}
\pending{与摘要、图表一致的主要结果及结论。}
\end{document}
